\documentclass[11pt]{article}

\usepackage[margin=1in]{geometry}
\usepackage{amsmath,amssymb,amsthm,mathtools}
\usepackage[hidelinks]{hyperref}

\newtheorem{theorem}{Theorem}[section]
\newtheorem{lemma}[theorem]{Lemma}
\newtheorem{proposition}[theorem]{Proposition}
\theoremstyle{remark}
\newtheorem{remark}[theorem]{Remark}

\newcommand{\tr}{\operatorname{tr}}
\newcommand{\Arg}{\operatorname{arg}}

\title{Positive Square Energy for a Two-$C_5$ Bouquet with Arbitrary Trees}
\author{}
\date{25 July 2026}

\begin{document}
\maketitle

\begin{abstract}
Let $G$ be a connected cactus whose only cyclic blocks are two copies of
$C_5$ having exactly one common vertex.  Arbitrary trees may be attached at
every vertex of this nine-vertex core.  We prove the explicit strict bound
\[
 s^+(G)\geq |V(G)|+1-\frac{4}{3\sqrt{13}}>|V(G)|.
\]
The proof eliminates the attached trees by exact matching belief propagation,
uses a coefficientwise integer-polynomial certificate on the resulting
weighted core, and integrates the phase of the normalized characteristic
polynomial.  The polynomial certificate is reproducible and is recorded by a
canonical SHA-256 digest.
\end{abstract}

\section{Introduction}

All graphs are finite, simple, and undirected.  If the adjacency eigenvalues
of $G$ are $\lambda_1,\ldots,\lambda_n$, put
\[
 s^+(G)=\sum_{\lambda_j>0}\lambda_j^2,
 \qquad
 s^-(G)=\sum_{\lambda_j<0}\lambda_j^2.
\]
Liu, Tang, and Zhang proved the positive and negative square-energy
conjecture for connected graphs~\cite{LTZ}.  Akbari, Kumar, Mohar, Pragada,
and Zhang proposed the stronger assertion $s^+(G)\geq n$ when a connected
$n$-vertex graph has at least $n+1$ edges~\cite{AKMPZ}.  Optimizing the
Liu--Tang--Zhang doubly nonnegative certificate block by block settles most
bicyclic cacti but leaves, in particular, the pair of cyclic block lengths
$\{5,5\}$.  The argument here resolves the shared-vertex member of that
family.

\begin{theorem}\label{thm:main}
Let $G$ be a connected cactus whose only cyclic blocks are two copies of
$C_5$, and suppose that these cycles share exactly one vertex $v_0$.
Every edge outside the cycles may belong to an arbitrary tree attached at any
vertex of the two-cycle core.  If $n=|V(G)|$, then
\begin{equation}\label{eq:main}
 s^+(G)\geq n+1-\frac{4}{3\sqrt{13}}>n.
\end{equation}
\end{theorem}

The hypothesis means a one-vertex bouquet: the two $C_5$ blocks intersect in
the single vertex $v_0$.  It does not cover vertex-disjoint cycles joined by a
bridge or a path, and it does not assert a result for arbitrary bicyclic
cacti.

The coefficient inequality used below is an equality when there are no
nontrivial tree messages.  The later analytic inequality
$\arctan u\leq u$ is strict for $u>0$, however, so we do not claim that the
constant in~\eqref{eq:main} is the best possible spectral lower bound.

\section{Signless matchings and grouped Sachs expansion}

For a forest $F$, define its signless matching partition by
\begin{equation}\label{eq:ZF}
 Z_F(t)=\sum_{M}t^{|V(F)|-2|M|},
\end{equation}
where the sum runs over all matchings of $F$.  More generally, for positive
vertex activities $a=(a_v)_{v\in V(H)}$, write
\begin{equation}\label{eq:weighted-Z}
 Z_H(a)=\sum_M\prod_{v\text{ unmatched by }M}a_v.
\end{equation}
All edge activities are one.  Thus~\eqref{eq:ZF} is the specialization of
\eqref{eq:weighted-Z} in which every vertex activity is $t$.

Let $\phi_G(z)=\det(zI-A(G))$ and normalize it on the imaginary axis by
\begin{equation}\label{eq:Psi}
 \Psi_G(t)=i^{-n}\phi_G(it)=\prod_{j=1}^n(t+i\lambda_j),
 \qquad t>0.
\end{equation}
Grouping the Sachs expansion by its collection $\mathcal C$ of
vertex-disjoint cycle components gives
\begin{equation}\label{eq:Sachs}
 \Psi_G(t)=
 \sum_{\mathcal C}
 \prod_{C\in\mathcal C}\bigl(-2i^{-|C|}\bigr)
 Z_{G-V(\mathcal C)}(t).
\end{equation}
The empty collection contributes $Z_G(t)$.  A $5$-cycle contributes the
multiplier $-2i^{-5}=2i$.  Since the two cycles under consideration share
$v_0$, they are not vertex-disjoint and cannot occur together in a Sachs
subgraph.  In particular, there is no double-cycle term in~\eqref{eq:Sachs}.

\section{Exact elimination of attached trees}

Orient every tree branch toward its core vertex.  For an oriented noncore
edge $u\to p$, let $T_{u\to p}$ be the component below $u$, including $u$,
and define
\begin{equation}\label{eq:q-definition}
 q_{u\to p}(t)=
 \frac{Z_{T_{u\to p}}(t)}{Z_{T_{u\to p}-u}(t)}.
\end{equation}
The denominator is a product of child-subtree partitions and is positive for
$t>0$.

\begin{lemma}[Rooted-tree matching recursion]\label{lem:BP}
For every oriented branch edge,
\begin{equation}\label{eq:BP}
 q_{u\to p}(t)=t+\sum_{w\text{ child of }u}\frac1{q_{w\to u}(t)}
 \geq t>0.
\end{equation}
After all noncore vertices are eliminated, the effective activity at a core
vertex $v$ is
\begin{equation}\label{eq:activity}
 a_v(t)=t+\sum_{u\text{ attached at }v}\frac1{q_{u\to v}(t)}
       =t+y_v(t),
 \qquad y_v(t)\geq0.
\end{equation}
Moreover, all terms of the grouped Sachs expansion have the same positive
prefactor
\begin{equation}\label{eq:K}
 K(t)=\prod_{u\text{ adjacent to the core}}Z_{T_{u\to v}}(t)>0.
\end{equation}
\end{lemma}

\begin{proof}
Split a matching of $T_{u\to p}$ according to whether $u$ is unmatched or is
matched to one of its children.  After the child-subtree partitions are
factored out, this gives~\eqref{eq:BP}.  A leaf has message $t$, so induction
also proves $q_{u\to p}\geq t>0$.

For a branch $T=T_{u\to v}$ rooted next to a core vertex $v$, a matching not
using $uv$ contributes $Z_T$.  A matching using $uv$ contributes $Z_{T-u}$.
Extracting $Z_T$ changes the latter choice into the weight
\[
 \frac{Z_{T-u}}{Z_T}=\frac1{q_{u\to v}}.
\]
The incident branch choices are mutually exclusive, so their sum adds exactly
the terms in~\eqref{eq:activity}.  Independence among branches yields the
factor~\eqref{eq:K}.

It remains important to check cycle deletion, rather than silently canceling
an unjustified ratio.  If a Sachs cycle deletes a core vertex $v$, every
branch formerly attached at $v$ becomes a disconnected full tree and
contributes $Z_T$, exactly its factor in $K(t)$.  At each retained core
vertex, the branch is absorbed into the same effective activity as before.
Thus deleting all five vertices of either cycle preserves the full common
factor $K(t)$.  No quotient of tree partitions remains.
\end{proof}

The same argument covers a rooted tree described as being identified with a
core vertex: remove that core root and treat each resulting component as one
of the branches above.

\section{The weighted two-$C_5$ core}

Label the common vertex $0$.  On each lobe $j\in\{1,2\}$, label the other
vertices consecutively $1,2,3,4$, with activities $x_1,x_2,x_3,x_4$ when one
lobe is under discussion.  Deleting vertex $0$ leaves a path $P_4$.  Its
weighted matching partition is
\begin{equation}\label{eq:A}
 A(x_1,x_2,x_3,x_4)
 =x_1x_2x_3x_4+x_3x_4+x_1x_4+x_1x_2+1.
\end{equation}
These are respectively the empty matching, the three one-edge matchings, and
the unique two-edge matching of $P_4$.

If vertex $0$ is matched along this lobe, the sum over the two possible
incident cycle edges is
\begin{equation}\label{eq:B}
 B(x_1,x_2,x_3,x_4)
 =x_2x_3x_4+x_2+x_4+x_1x_2x_3+x_1+x_3.
\end{equation}
Indeed, either edge from $0$ leaves a $P_3$.  The two displayed groups in
\eqref{eq:B} correspond to the two distinct choices of incident edge, so no
additional factor of two is present.

Let $A_j,B_j$ denote~\eqref{eq:A}--\eqref{eq:B} on lobe $j$, and let $a_0$
be the activity at the common vertex.  Splitting core matchings according to
the status of vertex $0$ gives
\begin{equation}\label{eq:R}
 R=a_0A_1A_2+B_1A_2+A_1B_2.
\end{equation}
The first term leaves $0$ unmatched; the other terms match it into exactly
one of the two lobes.  These cases are disjoint and exhaustive.

Deleting one whole $C_5$ from the core leaves the $P_4$ in the other lobe.
Lemma~\ref{lem:BP}, the $2i$ cycle multiplier, and the absence of a
double-cycle term therefore give the exact identity
\begin{equation}\label{eq:Psi-factor}
 \boxed{\ \Psi_G(t)=K(t)\bigl[R+2i(A_1+A_2)\bigr].\ }
\end{equation}
All quantities in brackets are positive for $t>0$.  Since $K(t)$ is positive
and real, it does not affect the continuous argument.  Consequently
\begin{equation}\label{eq:Theta-exact}
 \Theta_G(t):=\Arg\Psi_G(t)
 =\arctan\frac{2(A_1+A_2)}R,
 \qquad 0<\Theta_G(t)<\frac\pi2.
\end{equation}

\section{The coefficientwise certificate}

Put
\begin{equation}\label{eq:D}
 D(t)=t^4+7t^2+9.
\end{equation}
The algebraic heart of the proof is the following finite exact statement.

\begin{proposition}\label{prop:certificate}
Whenever all nine core activities are at least $t>0$,
\begin{equation}\label{eq:core-inequality}
 2R\geq tD(t)(A_1+A_2).
\end{equation}
Equality holds in~\eqref{eq:core-inequality} when every activity is $t$.
\end{proposition}

\begin{proof}
Write
\[
 a_0=t+y_0,
 \qquad x_{j,k}=t+y_{j,k}
 \quad(j=1,2,\ k=1,2,3,4),
\]
where every $y$ is nonnegative.  Substitute these expressions into
\eqref{eq:A}--\eqref{eq:R} and form
\begin{equation}\label{eq:P}
 P=2(a_0A_1A_2+B_1A_2+A_1B_2)
   -t(t^4+7t^2+9)(A_1+A_2).
\end{equation}
Appendix~\ref{app:certificate} gives an exact, exhaustive coefficient
verification over $\mathbb Z$.  The expansion has $1290$ nonzero monomials;
all coefficients are positive integers between $1$ and $22$, and every
monomial contains at least one $y$ variable.  Representative terms include
\[
 2t^8y_0,
 \qquad t^8y_{1,1},
 \qquad 2y_{2,1},
 \qquad 2y_{2,2}y_{2,3}y_{2,4}.
\]
Thus $P\geq0$ on the nonnegative orthant, proving
\eqref{eq:core-inequality}.  Since the part independent of all $y$ variables
is identically zero, equality holds when all $y$ vanish.
\end{proof}

Combining~\eqref{eq:Theta-exact} and~\eqref{eq:core-inequality}, and using
monotonicity of the arctangent, yields
\begin{equation}\label{eq:phase-bound}
 0<\Theta_G(t)
 \leq\arctan\frac{4}{tD(t)}.
\end{equation}
The factor $4$ is forced: the imaginary part in~\eqref{eq:Psi-factor} is
$2(A_1+A_2)$, while~\eqref{eq:core-inequality} bounds
$(A_1+A_2)/R$ by $2/(tD(t))$.

\section{Coulson integral and square energy}

For the continuous argument fixed in~\eqref{eq:Psi}, the signed
square-energy Coulson identity is
\begin{equation}\label{eq:Coulson}
 s^+(G)-s^-(G)
 =-\frac4\pi\int_0^\infty t\Theta_G(t)\,dt.
\end{equation}
For completeness, it follows by summing the scalar identity for the factors
$t+i\lambda_j$ and integrating; the trace relation
$\sum_j\lambda_j=0$ cancels the nonintegrable first-order term at infinity.
The first-quadrant location in~\eqref{eq:Theta-exact} fixes the branch and
precludes winding.

By~\eqref{eq:phase-bound} and $\arctan u\leq u$ for $u\geq0$,
\begin{equation}\label{eq:integral-bound}
 \int_0^\infty t\Theta_G(t)\,dt
 \leq\int_0^\infty\frac{4}{t^4+7t^2+9}\,dt.
\end{equation}
Set $t=\sqrt3u$.  The reciprocal-quartic formula
\[
 \int_0^\infty\frac{du}{u^4+bu^2+1}
 =\frac{\pi}{2\sqrt{2+b}}
 \qquad(b>-2)
\]
then gives the exact value
\begin{equation}\label{eq:rational-integral}
 \int_0^\infty\frac{4}{t^4+7t^2+9}\,dt
 =\frac{2\pi}{3\sqrt{13}}<\frac\pi2.
\end{equation}
Equations~\eqref{eq:Coulson}--\eqref{eq:rational-integral} imply
\begin{equation}\label{eq:difference}
 s^+(G)-s^-(G)\geq-\frac{8}{3\sqrt{13}}.
\end{equation}

The graph has exactly two independent cycles, hence $|E(G)|=n+1$.  Therefore
\begin{equation}\label{eq:sum}
 s^+(G)+s^-(G)=\tr A(G)^2=2|E(G)|=2n+2.
\end{equation}
Averaging~\eqref{eq:difference} and~\eqref{eq:sum} proves
\[
 s^+(G)\geq n+1-\frac{4}{3\sqrt{13}}>n,
\]
because $4<3\sqrt{13}$.  This proves Theorem~\ref{thm:main}.

\begin{remark}[Equality nuance]
The coefficientwise inequality~\eqref{eq:core-inequality} is exact at the bare
core, where every activity equals $t$.  Nevertheless,
$\arctan u<u$ for every $u>0$, and the phase in
\eqref{eq:Theta-exact} is positive.  Thus the displayed spectral bound is not
an equality even for the bare core, and no optimality claim for its constant
is made.
\end{remark}

\appendix
\section{Exact certificate and reproducibility}\label{app:certificate}

The certificate script is
\begin{center}
\texttt{positive-square-energy/experiments/\allowbreak
c5\_bouquet\_matching\_certificate.py}.
\end{center}
Run it from the repository root with
\begin{verbatim}
python positive-square-energy/experiments/c5_bouquet_matching_certificate.py
\end{verbatim}
It uses SymPy to construct~\eqref{eq:P} as a polynomial over $\mathbb Z$ in
the ordered variables
\[
 (t,y_0,y_{1,1},y_{1,2},y_{1,3},y_{1,4},
 y_{2,1},y_{2,2},y_{2,3},y_{2,4}).
\]
It then asserts all of the following:
\begin{enumerate}
\item the expanded polynomial has exactly $1290$ nonzero terms;
\item every coefficient is nonnegative (in fact the minimum is $1$ and the
      maximum is $22$);
\item substitution of zero for all nine $y$ variables gives zero;
\item every exponent vector has a nonzero $y$ coordinate; and
\item the SHA-256 digest of the canonical stream
      \texttt{monomial:coefficient}, one term per line in SymPy's ordered
      term order, is
\end{enumerate}
\begin{center}
\texttt{4c436cac772395d2a8edfdd81408ffe426759d3e94d66df2e4ab0235a3343110}.
\end{center}
The script prints
\begin{verbatim}
PASS two-C5 bouquet matching coefficient certificate
terms=1290 min_coefficient=1 max_coefficient=22
y_constant=0 all_coefficients_nonnegative=True
sha256=4c436cac772395d2a8edfdd81408ffe426759d3e94d66df2e4ab0235a3343110
\end{verbatim}
This is an exhaustive symbolic verification, not a numerical sampling of the
nonnegative orthant.  Since coefficientwise nonnegativity is itself the proof
of Proposition~\ref{prop:certificate}, the script is part of the exact proof.
For file-level provenance, the script file used for this manuscript has
SHA-256 digest
\begin{center}
\texttt{53e5ede12064953fe3d19aa173247e63ea3ab17d95ca2901d8f5707dae2e9f92}.
\end{center}

\begin{thebibliography}{9}

\bibitem{AKMPZ}
S.~Akbari, H.~Kumar, B.~Mohar, S.~Pragada, and S.~Zhang,
\emph{Refinement of a conjecture on positive square energy of graphs},
arXiv:2506.07264, 2025.

\bibitem{LTZ}
Y.~Liu, Q.~Tang, and S.~Zhang,
\emph{The positive and negative square-energy conjecture},
arXiv:2607.18031, 2026.

\end{thebibliography}

\end{document}
